% year: תשפ"ב
% type: מבחן
% semester: א
% moed: ב
% number: 4ב
% points: 9
% categories: ode
% source: exams/מבחנים/תשפב/מועד ב תשפב.pdf

\begin{question}
(9 נק') פתרו את המשוואה הדיפרנציאלית
\[ xe^y=(x^2+1)y' \]
\end{question}

\begin{hint}
זו משוואה פרידה: $e^{-y}y'=\frac{x}{x^2+1}$.
\end{hint}

\begin{solution}
כיוון ש-$x^2+1>0$ ו-$e^y>0$, נוכל לכתוב
\[ e^{-y}\,y'=\frac{x}{x^2+1}. \]
(אין פתרונות קבועים: עבור $y\equiv c$ נקבל $xe^c=0$ לכל $x$, וזה בלתי אפשרי.) נבצע אינטגרציה של שני האגפים:
\[ \int e^{-y}dy=\int\frac{x}{x^2+1}dx\ \Longrightarrow\ -e^{-y}=\frac12\ln(x^2+1)-C. \]
לכן $e^{-y}=C-\frac12\ln(x^2+1)$, ומכאן
\[ \boxed{y=-\ln\left(C-\tfrac12\ln(x^2+1)\right)},\qquad C\in\R, \]
בתחום שבו $C-\frac12\ln(x^2+1)>0$ (בפרט נדרש $C>0$; התחום הוא $|x|<\sqrt{e^{2C}-1}$).

בדיקה: $y'=\frac{\frac{x}{x^2+1}}{C-\frac12\ln(x^2+1)}=\frac{x}{x^2+1}e^{y}$, ולכן $(x^2+1)y'=xe^y$.
\end{solution}
